\documentclass[letterpaper,10pt]{article}

% !!! OVERLEAF COMPILER INSTRUCTION !!!
% To compile this Japanese document on Overleaf:
% 1. Click 'Menu' in the top left corner.
% 2. Change the 'Compiler' setting from 'pdfLaTeX' to 'XeLaTeX'.
% 3. Click 'Recompile'.

\usepackage{latexsym}
\usepackage[empty]{fullpage}
\usepackage{titlesec}
\usepackage{marvosym}
\usepackage[usenames,dvipsnames]{xcolor}
\usepackage{verbatim}
\usepackage{enumitem}

\definecolor{accent}{HTML}{000000} % Generic Black
\definecolor{creme}{HTML}{FCFBF7} % Subtle warm off-white
\pagecolor{creme}
\usepackage[colorlinks=true, urlcolor=accent, linkcolor=accent, pdfborder={0 0 0}]{hyperref}
\usepackage{fancyhdr}
\usepackage{tabularx}

% Japanese Support Package for XeLaTeX
\usepackage{xeCJK}
\setCJKmainfont[AutoFakeBold=true]{IPAexMincho} % Clean serif Japanese font
\setCJKsansfont[AutoFakeBold=true]{IPAexMincho}

% English font support
\usepackage[T1]{fontenc}
\usepackage{helvet}
\renewcommand{\familydefault}{\sfdefault}

\pagestyle{fancy}
\fancyhf{} % clear all header and footer fields
\fancyfoot{}
\renewcommand{\headrulewidth}{0pt}
\renewcommand{\footrulewidth}{0pt}

% Adjust margins for 1-page fit
\addtolength{\oddsidemargin}{-0.55in}
\addtolength{\evensidemargin}{-0.55in}
\addtolength{\textwidth}{1.1in}
\addtolength{\topmargin}{-0.55in}
\addtolength{\textheight}{1.15in}

\urlstyle{same}

\raggedbottom
\raggedright
\setlength{\tabcolsep}{0in}

% Sections formatting
\titleformat{\section}{
  \vspace{0pt}\bfseries\raggedright\Large\color{accent}
}{}{0em}{}[\color{accent}\titlerule \vspace{4pt}]

% Custom commands
\newcommand{\resumeItem}[1]{
  \item\small{#1}
}

\newcommand{\resumeSubheading}[4]{
  \item
    \begin{tabular*}{0.98\textwidth}[t]{l@{\extracolsep{\fill}}r}
      \textbf{#1} & \small#2 \\
      \textit{\small#3} & \textit{\small #4} \\
    \end{tabular*}\vspace{2pt}
}

\newcommand{\resumeSubSubheading}[2]{
    \item
    \begin{tabular*}{0.98\textwidth}[t]{l@{\extracolsep{\fill}}r}
      \textit{\small#1} & \textit{\small #2} \\
    \end{tabular*}\vspace{2pt}
}

\newcommand{\resumeProjectHeading}[2]{
    \item
    \begin{tabular*}{0.98\textwidth}[t]{l@{\extracolsep{\fill}}r}
      \small#1 & \small#2 \\
    \end{tabular*}\vspace{2pt}
}

\newcommand{\resumeSubItem}[1]{\resumeItem{#1}}

\renewcommand\labelitemii{$\vcenter{\hbox{\tiny$\bullet$}}$}

\newcommand{\resumeSubHeadingListStart}{\begin{itemize}[leftmargin=0.12in, label={}, itemsep=5pt, parsep=0pt, topsep=2pt]}
\newcommand{\resumeSubHeadingListEnd}{\end{itemize}}
\newcommand{\resumeItemListStart}{\begin{itemize}[leftmargin=0.18in, itemsep=3pt, parsep=0pt, topsep=1pt, partopsep=0pt]}
\newcommand{\resumeItemListEnd}{\end{itemize}\vspace{1pt}}

\begin{document}

\begin{center}
    \textbf{\Huge \scshape インドラニール・サマンタ} \\ \vspace{6pt}
    \small \href{https://linkedin.com/in/indraneel-samanta}{\underline{LinkedIn}} $|$
    \href{https://github.com/NotCatfish}{\underline{GitHub}} $|$
    \href{https://indraneelsamanta.vercel.app/}{\underline{Portfolio}}
\end{center}

%-----------SUMMARY-----------
\section{目的}
  \small{知的好奇心とAI等の新技術への情熱を持つソフトウェアエンジニア志望。大胆な挑戦を恐れず、他者が見過ごす斬新なアプローチに挑み、革新的かつユーザー本位のプロダクトを共創したい。スケーラブルな機械学習システムとシームレスなフルスタックUXの架け橋となることに尽力する。}
\vspace{4pt}

%-----------EDUCATION-----------
\section{学歴}
  \resumeSubHeadingListStart
    \resumeSubheading
      {Dwarkadas J. Sanghvi 工科大学 (DJSCE)}{インド、ムンバイ}
      {人工知能・機械学習（AIML）学士課程}{2024 -- 2028 (卒業予定)}
  \resumeSubHeadingListEnd

%-----------SKILLS & LANGUAGES-----------
\section{スキルと言語}
 \begin{itemize}[leftmargin=0.15in, label={}, itemsep=4pt, parsep=0pt, topsep=2pt]
    \small{\item{
     \textbf{プログラミングとデータベース}{: C, C++, Python, SQLite, PostgreSQL} \\[3pt]
     \textbf{データとML}{: Pandas, Plotly, Jupyter, Scikit-Learn, XGBoost, EDA, Predictive Modeling, AI Co-creation} \\[3pt]
     \textbf{Webとクラウド}{: Next.js, Supabase, Git} \\[3pt]
     \textbf{使用言語}{: 英語 (流暢), 日本語 (N1準備中), ヒンディー語 (流暢), ベンガル語 (流暢), グジャラート語 (リスニング理解), マラーティー語 (リスニング理解)}
    }}
 \end{itemize}

%-----------PROJECTS-----------
\section{テクニカルプロジェクト}
    \resumeSubHeadingListStart
      \resumeProjectHeading
          {\href{https://github.com/NotCatfish/Spotify-Analytics-Pipeline}{\textbf{Spotify MLパイプラインとデータエンジニアリング}} | \emph{Python, Pandas, PostgreSQL, XGBoost, Scikit-Learn}}{}
          \resumeItemListStart
            \resumeItem{\textbf{メモリ使用量を79\%（277MBから58MBへ）削減:} 生のストリーミングログをデータベースに圧縮するデータパイプラインを構築。}
            \resumeItem{\textbf{60以上のインタラクティブチャートを自動生成するツールを作成:} 手動設定なしで視聴習慣を即座に表示。}
            \resumeItem{\textbf{短期的なユーザー習慣を追跡:} スキップ率が\textbf{31\%}から\textbf{4\%}へ急減した際、システムが迅速に適応できるよう支援。}
            \resumeItem{\textbf{スキップを特定する予測モデルを訓練:} 単なる推測を避けるため、不均衡データを処理するようシステムを調整。}
            \resumeItem{\textbf{スキップ曲検出モデルの能力向上:} データ漏洩を防ぐため、将来のタイムラインで厳密にテスト。}
            \resumeItem{\textbf{隠れたリスニングパターンを抽出:} 生のタイムスタンプから\textbf{40}以上のデータポイントを作成（時間帯ごとの傾向など）。}
          \resumeItemListEnd

      \resumeProjectHeading
          {\href{https://github.com/NotCatfish/Otakufy}{\textbf{Otakufy フルスタック日本語学習プラットフォーム}} | \emph{Next.js, React, Tailwind CSS, Supabase, Vercel}}{Live: \href{https://otakufy.vercel.app/}{\underline{otakufy.vercel.app}}}
          \resumeItemListStart
            \resumeItem{\textbf{全JLPTレベルで超高速クイズ読込を実現:} 語彙データを動的取得せず事前読込するWebアプリを構築。}
            \resumeItem{\textbf{モバイルライクな高応答性UXを構築:} カスタムフラッシュカードエンジンをフロントエンドに統合。}
            \resumeItem{\textbf{ユーザーエンゲージメントを促進:} 安全なログインと進捗追跡により、連続学習とリーダーボードを提供。}
            \resumeItem{\textbf{Webサイト全体をゼロから迅速に構築:} 明確な指示に基づきAIコーディングアシスタントを活用。}
            \resumeItem{\textbf{全デバイスでシームレスな体験を保証:} タブレットやPCに自動拡張するモバイルファーストなレイアウトを設計。}
          \resumeItemListEnd

      \resumeProjectHeading
          {\href{https://github.com/NotCatfish/portfolio}{\textbf{個人ポートフォリオ ＆ インタラクティブ履歴書}} | \emph{React 19, Tailwind CSS v4, Vite, Framer Motion, i18n}}{Live: \href{https://indraneelsamanta.vercel.app/}{\underline{indraneelsamanta.vercel.app}}}
          \resumeItemListStart
            \resumeItem{\textbf{高速でレスポンシブな個人サイトを構築:} スムーズなスクロールアニメーションを備えたモダンなツールを使用。}
            \resumeItem{\textbf{日英翻訳スイッチとカスタムダークモードを追加:} 日本の採用担当者がサイトを読みやすいように設計。}
            \resumeItem{\textbf{実用的なお問い合わせフォームを設置:} スパム対策機能を組み込みつつ、基本コードをクリーンに維持。}
            \resumeItem{\textbf{Webサイト公開プロセスを自動化:} 手作業なしで即座かつシームレスにライブサイトにプッシュされるように設定。}
          \resumeItemListEnd
    \resumeSubHeadingListEnd

%-----------EXPERIENCE-----------
\section{職歴}
  \resumeSubHeadingListStart
    \resumeProjectHeading
      {\textbf{ロジスティクス・イベント共同委員} | \emph{Google Developer Student Club (GDSC) DJSCE}}{2025年10月 -- 現在}
      \resumeItemListStart
        \resumeItem{\textbf{5名のチームを率いて、350名以上の参加者のためのロジスティクス、環境づくり、イベント後の運営を調整}し、リーダーシップを発揮。}
        \resumeItem{\textbf{スポンサーや参加者にアプローチ}し、メールや直接の電話連絡を通じてイベントのサインアップを促進。}
      \resumeItemListEnd
  \resumeSubHeadingListEnd

\end{document}
